\documentclass[11pt,a4paper]{article}

%====================================================
% PACKAGES
%====================================================
\usepackage[utf8]{inputenc}
\usepackage[T1]{fontenc}
\usepackage[french]{babel}
\usepackage[a4paper,margin=2cm]{geometry}
\usepackage{amsmath,amssymb,amsthm,mathtools}
\usepackage{enumitem}
\usepackage{fancyhdr}
\usepackage{lastpage}
\usepackage{xcolor}
\usepackage{array}
\usepackage{booktabs}

%====================================================
% MISE EN PAGE
%====================================================
\setlength{\headheight}{14pt}
\pagestyle{fancy}
\fancyhf{}
\fancyhead[L]{Classe préparatoire}
\fancyhead[C]{Devoir surveillé -- Algèbre linéaire}
\fancyhead[R]{Page \thepage/\pageref{LastPage}}
\fancyfoot[C]{}

\setlength{\parindent}{0pt}
\setlength{\parskip}{0.4em}

\definecolor{bleu}{RGB}{20,60,120}
\definecolor{gris}{RGB}{85,85,85}

\newcommand{\pts}[1]{\hfill \textbf{[#1 pts]}}
\newcommand{\R}{\mathbb{R}}
\newcommand{\Ker}{\operatorname{Ker}}
\newcommand{\Img}{\operatorname{Im}}
\newcommand{\rg}{\operatorname{rg}}
\newcommand{\Mat}{\operatorname{Mat}}
\newcommand{\Vect}{\operatorname{Vect}}
\newcommand{\Id}{\operatorname{Id}}

\newcounter{probleme}
\newenvironment{probleme}[2]{
    \refstepcounter{probleme}
    \vspace{0.5cm}
    {\Large\color{bleu}\textbf{Problème \theprobleme{} -- #1}}\hfill \textbf{#2 pts}
    \par\medskip
    \hrule
    \medskip
}{\vspace{0.4cm}}

\newcommand{\partie}[1]{
    \vspace{0.25cm}
    {\large\color{gris}\textbf{#1}}
    \par\smallskip
}

\setlist[enumerate,1]{label=\textbf{\arabic*.}, leftmargin=1.1cm}
\setlist[enumerate,2]{label=\textbf{\alph*)}, leftmargin=1cm}
\setlist[itemize]{leftmargin=1cm}

%====================================================
% DOCUMENT
%====================================================
\begin{document}

\begin{center}
    {\Huge \textbf{Devoir surveillé}}\\[0.2cm]
    {\Large \textbf{Algèbre linéaire}}\\[0.2cm]
    {\large Systèmes linéaires, matrices, applications linéaires, rang et déterminants}\\[0.4cm]
\end{center}

\begin{center}
\begin{tabular}{|>{\bfseries}c|c|}
\hline
Durée indicative & 3h \\
\hline
Niveau & Bac + 1 -- classe préparatoire scientifique \\
\hline
Total & 85 points \\
\hline
\end{tabular}
\end{center}

\vspace{0.3cm}

\textbf{Consignes.}
\begin{itemize}
    \item La qualité de la rédaction, la clarté des raisonnements et la justification des réponses seront prises en compte.
    \item Les calculs doivent être menés proprement. Toute réponse non justifiée peut ne pas recevoir la totalité des points.
    \item Les trois problèmes sont indépendants.
    \item Aucune notion de réduction, valeur propre, diagonalisation, polynôme d'endomorphisme ou produit scalaire n'est attendue.
    \item Les dernières questions de chaque problème sont plus ouvertes et demandent davantage de recul.
\end{itemize}

\vspace{0.2cm}

\begin{center}
\begin{tabular}{|c|c|c|}
\hline
Problème & Thème principal & Barème \\
\hline
1 & Systèmes linéaires, matrices, rang et compatibilité & 28 pts \\
\hline
2 & Machine linéaire, inverse et pertes d'information & 27 pts \\
\hline
3 & Noyau, image, rang, déterminant et généralisation & 30 pts \\
\hline
\textbf{Total} & & \textbf{85 pts} \\
\hline
\end{tabular}
\end{center}

%====================================================
% PROBLEME 1
%====================================================
\begin{probleme}{Un système linéaire à paramètre}{28}

On considère, pour un réel \(m\), le système linéaire \((S_m)\) d'inconnues \(x,y,z\) :
\[
(S_m)\qquad
\begin{cases}
x+y+z=1,\\
x+my+z=2,\\
x+y+mz=3.
\end{cases}
\]

On note
\[
A_m=
\begin{pmatrix}
1&1&1\\
1&m&1\\
1&1&m
\end{pmatrix},
\qquad
X=
\begin{pmatrix}
x\\y\\z
\end{pmatrix},
\qquad
B=
\begin{pmatrix}
1\\2\\3
\end{pmatrix}.
\]

\partie{Partie A -- Étude classique du système}

\begin{enumerate}
    \item Écrire le système \((S_m)\) sous la forme matricielle
    \[
    A_mX=B.
    \]
    \pts{1}

    \item Calculer le déterminant de \(A_m\). \pts{3}

    \item Pour quelles valeurs de \(m\) la matrice \(A_m\) est-elle inversible ? \pts{2}

    \item On suppose dans cette question que \(m\neq 1\).
    Résoudre le système \((S_m)\). \pts{5}

    \item Étudier séparément le cas \(m=1\).
    Le système est-il compatible ? Si oui, donner l'ensemble de ses solutions. \pts{3}

    \item En déduire, selon les valeurs de \(m\), si le système \((S_m)\) admet :
    \begin{itemize}
        \item une unique solution ;
        \item aucune solution ;
        \item une infinité de solutions.
    \end{itemize}
    \pts{2}
\end{enumerate}

\partie{Partie B -- Rang et système homogène associé}

On considère le système homogène associé :
\[
(H_m)\qquad A_mX=0.
\]

\begin{enumerate}[resume]
    \item Déterminer le rang de \(A_m\) selon les valeurs de \(m\). \pts{3}

    \item Résoudre le système homogène \((H_m)\) selon les valeurs de \(m\). \pts{3}

    \item Expliquer pourquoi le résultat obtenu est cohérent avec le théorème du rang appliqué à l'application linéaire
    \[
    u_m:\R^3\longrightarrow\R^3,
    \qquad
    X\longmapsto A_mX.
    \]
    \pts{3}
\end{enumerate}

\partie{Partie C -- Généralisation}

On remplace maintenant le second membre \(B\) par un second membre quelconque
\[
Y=
\begin{pmatrix}
a\\b\\c
\end{pmatrix}
\in\R^3.
\]
On considère donc le système
\[
A_mX=Y.
\]

\begin{enumerate}[resume]
    \item Lorsque \(m\neq 1\), justifier que le système admet une unique solution pour tout \(Y\in\R^3\). \pts{2}

    \item On suppose maintenant \(m=1\). Donner une condition nécessaire et suffisante portant sur \(a,b,c\) pour que le système
    \[
    A_1X=Y
    \]
    soit compatible. \pts{3}

    \item Dans le cas \(m=1\), interpréter géométriquement l'ensemble des seconds membres \(Y\) pour lesquels le système est compatible. \pts{1}
\end{enumerate}

\end{probleme}

%====================================================
% PROBLEME 2
%====================================================
\begin{probleme}{La machine à mélanger les jetons}{27}

Une petite machine reçoit trois nombres réels \(x,y,z\), représentant les quantités de trois types de jetons : rouges, bleus et verts. Elle renvoie trois nouvelles quantités \(a,b,c\) selon la règle :
\[
\begin{cases}
a=x+y,\\
b=y+z,\\
c=x+z.
\end{cases}
\]

On définit ainsi une application
\[
T:\R^3\longrightarrow \R^3,
\qquad
T(x,y,z)=(x+y,\ y+z,\ x+z).
\]

\partie{Partie A -- Étude de la machine complète}

\begin{enumerate}
    \item Montrer que \(T\) est une application linéaire. \pts{2}

    \item Donner la matrice \(M\) de \(T\) dans la base canonique de \(\R^3\). \pts{2}

    \item Calculer le déterminant de \(M\). \pts{3}

    \item En déduire que \(T\) est bijective. \pts{2}

    \item La machine affiche en sortie
    \[
    (a,b,c)=(7,8,9).
    \]
    Déterminer les valeurs initiales \(x,y,z\). \pts{3}

    \item Résoudre plus généralement le système
    \[
    T(x,y,z)=(a,b,c),
    \]
    où \(a,b,c\) sont trois réels quelconques. \pts{4}

    \item En déduire la matrice \(M^{-1}\). \pts{3}
\end{enumerate}

\partie{Partie B -- Machine avec perte d'information}

On modifie maintenant la machine. Elle renvoie seulement les deux nombres
\[
u=x+y,
\qquad
v=y+z.
\]
On définit donc
\[
F:\R^3\longrightarrow \R^2,
\qquad
F(x,y,z)=(x+y,\ y+z).
\]

\begin{enumerate}[resume]
    \item Donner la matrice \(N\) de \(F\) dans les bases canoniques. \pts{1}

    \item Déterminer \(\Ker(F)\). \pts{3}

    \item En déduire le rang de \(F\) à l'aide du théorème du rang. \pts{2}

    \item L'application \(F\) est-elle injective ? Est-elle surjective ? Justifier. \pts{3}
\end{enumerate}

\partie{Partie C -- Une généralisation à paramètre}

Pour un réel \(\lambda\), on définit une nouvelle machine
\[
T_\lambda:\R^3\longrightarrow\R^3
\]
par
\[
T_\lambda(x,y,z)=(x+y,\ y+z,\ x+\lambda z).
\]

\begin{enumerate}[resume]
    \item Donner la matrice \(M_\lambda\) de \(T_\lambda\) dans la base canonique. \pts{1}

    \item Calculer \(\det(M_\lambda)\). \pts{2}

    \item Pour quelles valeurs de \(\lambda\) la machine \(T_\lambda\) permet-elle toujours de retrouver les valeurs initiales \(x,y,z\) à partir de la sortie ? \pts{2}
\end{enumerate}

\end{probleme}

%====================================================
% PROBLEME 3
%====================================================
\begin{probleme}{Une application linéaire presque symétrique}{30}

On considère l'application
\[
u:\R^3\longrightarrow \R^3
\]
définie par
\[
u(x,y,z)=(x+y,\ y+z,\ x+z).
\]

On note \(B\) la base canonique de \(\R^3\).

\partie{Partie A -- Une application bijective}

\begin{enumerate}
    \item Justifier que \(u\) est linéaire. \pts{1}

    \item Donner la matrice \(A=\Mat_B(u)\). \pts{2}

    \item Calculer \(\det(A)\). \pts{3}

    \item En déduire que \(u\) est un automorphisme de \(\R^3\). \pts{2}

    \item Calculer explicitement \(\Ker(u)\). Vérifier que le résultat est cohérent avec la question précédente. \pts{3}

    \item Déterminer \(\Img(u)\). \pts{2}
\end{enumerate}

\partie{Partie B -- Une application proche mais non bijective}

On considère maintenant l'application
\[
v:\R^3\longrightarrow \R^3
\]
définie par
\[
v(x,y,z)=(x+y,\ y+z,\ x+2y+z).
\]

\begin{enumerate}[resume]
    \item Montrer que \(v\) est linéaire et donner sa matrice \(C\) dans la base canonique. \pts{2}

    \item Calculer \(\det(C)\). \pts{2}

    \item Déterminer \(\Ker(v)\). \pts{3}

    \item En déduire le rang de \(v\) à l'aide du théorème du rang. \pts{2}

    \item Donner une base de \(\Img(v)\). \pts{3}

    \item Expliquer clairement pourquoi \(v\) n'est pas bijective. \pts{2}
\end{enumerate}

\partie{Partie C -- Une généralisation élégante}

Pour un réel \(\alpha\), on définit
\[
w_\alpha:\R^3\longrightarrow\R^3
\]
par
\[
w_\alpha(x,y,z)=(x+y,\ y+z,\ x+\alpha y+z).
\]

\begin{enumerate}[resume]
    \item Donner la matrice \(D_\alpha\) de \(w_\alpha\) dans la base canonique. \pts{1}

    \item Calculer \(\det(D_\alpha)\). \pts{3}

    \item Pour quelles valeurs de \(\alpha\) l'application \(w_\alpha\) est-elle un automorphisme de \(\R^3\) ? \pts{2}

    \item Dans le cas où \(w_\alpha\) n'est pas un automorphisme, déterminer \(\Ker(w_\alpha)\). \pts{3}

    \item Dans ce même cas, donner le rang de \(w_\alpha\), puis une base de \(\Img(w_\alpha)\). \pts{3}

    \item Comparer les applications \(u\), \(v\) et \(w_\alpha\) du point de vue de l'inversibilité, du noyau et du rang. Rédiger une conclusion synthétique. \pts{2}
\end{enumerate}

\end{probleme}

\vfill

\begin{center}
    \textbf{Fin du sujet}
\end{center}

\end{document}